% !TeX root = ../Book1.tex
\begin{figure}[htbp]
	\centering
	\begin{tikzpicture}[scale=0.95,font=\small]
		\begin{scope}[xshift=-3.1cm]
			\fill[GeometryBlueLight] (0,0.9) rectangle (1.7,3.1);
			\draw[->] (-0.7,0)--(2.5,0) node[right] {\(x\)};
			\draw[GeometryBlue,thick] (0,0.35)--(0,3.25);
			\draw[GeometryBlue,thick] (1.7,0.35)--(1.7,3.25);
			\draw[GeometryBlue,dashed] (-0.4,0.9)--(2.1,0.9);
			\node[left] at (-0.4,0.9) {\(y=Y\)};
			\draw[->,GeometryBlue] (0,1.5)--(1.7,1.5) node[midway,above] {\(z\mapsto z+1\)};
			\draw[->,GeometryBlue] (0.85,2.2)--(0.85,3.2);
			\node[above] at (0.85,3.2) {尖端 \(y\to\infty\)};
			\node[below] at (0,0) {\(0\)};
			\node[below] at (1.7,0) {\(1\)};
			\node[align=center] at (0.85,-0.85) {平移商的尖点端\\竖边配对，阴影部分面积有限};
		\end{scope}
		\begin{scope}[xshift=3.7cm]
			\fill[GeometryBlueLight] (2.1,0) arc (0:180:2.1)--(-1.05,0) arc (180:0:1.05)--cycle;
			\draw[->] (-2.5,0)--(2.5,0) node[right] {\(x\)};
			\draw[GeometryBlue,thick] (2.1,0) arc (0:180:2.1);
			\draw[GeometryBlue,thick] (1.05,0) arc (0:180:1.05);
			\draw[GeometryGray,dashed] (0,0)--(0,2.7);
			\draw[GeometryBlue,very thick] (0,1.05)--(0,2.1);
			\node[right] at (0.06,1.6) {轴};
			\draw[->,GeometryBlue] (140:1.05)--(140:2.1);
			\node[above,align=center] at (-1.15,2.4) {\(z\mapsto e^\ell z\)};
			\node at (-1.5,0.45) {\(r<0\)};
			\node at (1.5,0.45) {\(r>0\)};
			\node[align=center] at (0,-0.85) {缩放商的基本区域\\两条半圆配对，轴两侧为漏斗};
		\end{scope}
	\end{tikzpicture}
	\caption[尖点与漏斗的基本区域]{左图截取 \(0\le x\le1,y\ge Y\)，仅表示平移商中的尖点端；右图取 \(1\le|z|\le e^\ell\)，表示整个缩放圆柱的一个基本区域。右图以半径比 \(2\) 示意，粗轴段配对后成为闭测地线。图中的 Euclid 大小不表示双曲长度。}
	\label{g1:fig:cusp-funnel-domains}
\end{figure}
